\subsection{The coefficient lattice of the fitted test}
\label{subsec:fitted-test-slope}

At the other end of the height comparison the coefficient functionals must act on an integral lattice. Its highest-weight space and dual fix the normalization of each grid test.

\begin{established}[Rank and metric of the fitted dual lattice]\label{hyp:fitted-dual-lattice}
For the fitted top test let \(r=k+m-1\) and \(H\) be the Hermitian
Gram matrix on a definite class lattice in its chosen frame.
The rank and the Gram matrix of the test on the dual are
\[
 n_*=(k+m)+m(k+m+4)=(r+1)+m(r+5),\qquad
 J=(H^{-1})^{\mathsf t},\quad \det J=(\det H)^{-1}.
\]
The target of the Fourier coefficient map is the bilinear dual of
the orthogonal sum of the transported differential Weyl lattices over
all the test data and classes.
\end{established}
\begin{proof}
The base contributes \(k+m\) rows and every one of the \(m\) bands
contributes \(k+m+4\). At each finite place the class representative
acts isometrically. For its local lattice \(L_v\),
\[
              (g_v L_v)^\#=g_v L_v^\#=g_v L_v .
\]
Intersecting the local lattices in the rational space gives the same
identity globally. At infinity the definite form gives a positive
matrix in every selected coefficient embedding, as the embeddings
respect involution. Evaluation pairs a frame with its bilinear dual.
Transporting the form through this pairing gives the inverse
transpose on each summand, and taking determinants proves the formula.
\end{proof}

\begin{proposition}[Maximal slope of the fitted coefficient lattice]\label{prop:fitted-coefficient-slope}
In the test dual \(M\) of Lemma~\ref{hyp:fitted-dual-lattice}, let \(R_*\) be the
weight-spread bound on the selected fitted grid. Uniformly in the blocks and twists in the prescribed parameter ranges,
\[
 \mu_{\max}(M)\le4fn_*R_*\log(n_*R_*).
\]

Here \(n_*\le2mr\), \(R_*\le CmN^3b\), and
\(\log(n_*R_*)\le4\log b\).
\end{proposition}

\begin{proof}
\proofstep{Compare the integral Weyl lattice with the column-symmetric tensor lattice.}
Let \(n\) be the
underlying rank and \(R\ge1\) a bound for the spread of
each dominant weight \(\lambda_\tau\) at the \(f\) labels.
For each label remove the determinant twist
\(\lambda_{\tau,n}\), leaving a partition
\(\nu_\tau\) whose first part is at most \(\lfloor R\rfloor\).
The standard highest vector is the product of leading minors
with multiplicities \(\nu_{\tau,i}-\nu_{\tau,i+1}\).

For a complex polynomial \(P\) in the matrix variables,
the unnormalized Fischer inner product has squared norm
\(\sum_\alpha\alpha!|P_\alpha|^2\).
Let \(s_\tau\) be the squared norm of this highest vector.
The sum of the absolute coefficients is at most
\((n!)^{\lfloor R\rfloor}\), by expanding all the minors.
Each monomial has degree \(\nu_{\tau,j}\) in column \(j\),
so 
\[
 \prod_i\alpha_{ij}!=\nu_{\tau,j}!\binom{\nu_{\tau,j}}{\alpha_{1j},\ldots,\alpha_{nj}}^{-1}\le\nu_{\tau,j}!\le\lfloor R\rfloor!.
\]
In particular \(s_\tau\) is a positive integer and
\[
  s_\tau\le (n!)^{2\lfloor R\rfloor}
                            (\lfloor R\rfloor!)^n.
\]
The tensor pairing on the representation is normalized by
\(s=\prod_\tau s_\tau\). Thus
\[
                 \log s\le 4fnR\log(nR).
\]

The integral comparison uses the two standard Weyl lattices
in the polynomial realization. For the Fischer pairing
normalized by \(s\), the lattice obtained by divided powers
of the row lowering operators on the extremal vector and
the lattice spanned by the integral bideterminants are mutual
duals. In the tableau proof the row words pair with the
semistandard bideterminants by the coefficients of lower
unipotent substitution. The leading tableau monomials give
a triangular family with diagonal entries one. This yields
dual bases over the integers for the normalized form and the
same duality over each coefficient ring.
The duality is computed at every finite place in paired
frames. More explicitly, at a coefficient prime \(w\) choose
a frame at \(\overline w\) and its Hermitian dual frame at
\(w\). For the transported basis \(b_i\) and its Hermitian dual \(b_j^\vee\), the transition Gram has entries
\[
  h(b_i,b_j^\vee)=\delta_{ij}.
\]

In those coordinates the pairing of the two
standard lattices is perfect. Thus the Hermitian dual of
the transported differential Weyl lattice is the transported
bideterminant lattice of the Hermitian dual
underlying frame. In the self-dual case these two
underlying frames are bases of the same underlying lattice
at the indicated places. For two vectors in the transported
bideterminant lattice at \(\overline w\) and \(w\), their
transition Gram has integral entries and unit determinant.
After multiplication by \(s\), the pairing is the integer
coefficient Fischer pairing with factorial weights,
evaluated after a substitution by this integral Gram.
The determinant twist contributes a unit even if
\(\lambda_{\tau,n}\) is negative.
For a transition Gram \(G\) between the two integral frames this last assertion reads
\[
  v_w(\det G)=0\quad\Longrightarrow\quad v_w\big((\det G)^{\lambda_{\tau,n}}\big)=0\qquad(\lambda_{\tau,n}\in\mathbb Z).
\]
 This proves integrality
at all \(w\). Summing over the orthogonal class blocks
preserves integrality with the same \(s\).
\label{cor:fitted-test-slope-1}

Let \(H\) denote the rational Gram of that sum and \(L\)
its test lattice. For \(x\) in its bilinear dual the vector
\(H^{-1}\overline{x}\) lies in the Hermitian dual.
Pairing the images of \(y,x\) in that order gives
\(\overline{x}^{\,t}(H^{-1})^t y\). Hence the scaled
target form \(s(H^{-1})^t\) is integral on \(M\).
For the determinant comparison take a rank-\(j\)
sublattice with rational basis matrix \(U\) and determinant
ideal \(I\) in that basis. Write \(J=(H^{-1})^t\), let
\(c\) be the coefficient involution and put
\(\Delta=\det(c(U)^tJU)\). On the determinant line the two conclusions are
\[
       s^j\Delta\,c(I)I\subseteq O_K,\qquad \Delta\ne0 .
\]
The inclusion follows from integrality of the scaled form.

For the
second assertion, \(c\) is complex conjugation at the selected
embedding and the Gram form there is positive definite.
For each embedding \(\sigma\) put \(\sigma'=\overline{\sigma\circ c}\)
and write \(G_\sigma=\overline{\sigma(U)}^{\,t}\sigma(J)\sigma(U)\).
The entries of \(J\) lie in the CM subfield on which \(c\) is
conjugation, so \(\sigma'(J)=\sigma(J)\). Taking the exterior
power of the mixed Gram pairing therefore gives
\[
 |\sigma(\Delta)|^2
       \le\det G_{\sigma'}\,\det G_\sigma,\qquad
 \frac{\log|\mathrm N_{K/\mathbb Q}\Delta|}{[K:\mathbb Q]}
       \le \frac{1}{[K:\mathbb Q]}\sum_\sigma\log\det G_\sigma .
\]
The first bound is Cauchy--Schwarz for the two exterior basis
vectors. Summing its logarithm gives the second because
\(\sigma\mapsto\sigma'\) permutes the embeddings.
Taking norms of the integral ideal gives
\[
 -j\log s\le
  \frac{2\log\mathrm N I}{[K:\Q]}
  +\operatorname{mean}_{\sigma}\log
                     \det(\sigma(U)^*\sigma(J)\sigma(U)).
\]
The right side is the negative degree of the sublattice.
This proves \(\mu_{\max}(M)\le\log s\) regardless of the
number of test blocks.
\label{cor:fitted-test-slope-2}

Write \(L_c\) for the maximum of the nonnegative differences
\(v_2-v_k\) between the head weights at each index, and
let \(g(C)\) be the gap parameter of the polynomial size
bounds. \proofstep{Bound the spread of the actual fitted weights.}
For the actual weight spread use
\[
       R_*=(m+1)(g(C)+1)+L_c+n_*.
\]
The top test uses comparison parameter one in its highest
weight. Before its index shifts, the ordered test entries
lie between \(v_k-1/2\) and
\(v_2+1/2+(m+1)g(1)\). The lower endpoint bounds the
smallest head weight. The band and top entries lie
between that endpoint and the top endpoint by the
interlacing inequalities. The conversion from the ordered
entries to the highest weight adds an index difference
of at most \(n_*\).

Since \(g(1)\le g(C)\) for \(C\ge1\),
this proves the bound \(R_*\) on every weight difference.
Therefore the fitted target satisfies
\[
             \mu_{\max}(M)\le4fn_*R_*\log(n_*R_*).
\]
The selected size conditions give \(n_*\le2mr\),
\(R_*\le CmN^3b\) and \(\log(n_*R_*)\le4\log b\).
The bound is uniform over the test blocks and the twists
at every point of this size record.
\label{cor:fitted-test-slope-3}
\end{proof}

